\documentclass[preprint,12pt]{elsarticle}
\usepackage[T1]{fontenc}
\usepackage[utf8]{inputenc}
\usepackage[brazil]{babel}
\usepackage{amsmath,amssymb,booktabs,graphicx,xurl,placeins}
\usepackage[hidelinks]{hyperref}
\setlength{\emergencystretch}{2em}
\renewcommand{\floatpagefraction}{0.8}
\renewcommand{\topfraction}{0.9}
\renewcommand{\textfraction}{0.1}
\journal{Supply Chain Analytics -- manuscrito de trabalho em português}
\begin{document}
\begin{frontmatter}
\title{AIPE: uma metodologia modular para seleção contextual de políticas de reposição}
\author[ufal]{Matheus Inácio Batista Santos}
\affiliation[ufal]{organization={Instituto de Computação, Universidade Federal de Alagoas},country={Brasil}}
\begin{abstract}
Gerentes de logística e estoques precisam conciliar custo e atendimento sob demandas heterogêneas. Propomos o Adaptive Inventory Policy Engine (AIPE), uma metodologia modular de apoio à decisão que vincula contexto, evidência operacional, seleção e avaliação. Seu Policy Selection Engine (PSE) associa perfis a famílias de políticas ou estima recomendações por série, com uma regra explícita para substituir ou conservar a referência. Avaliamos sete candidatas em duas aplicações separadas. Na Bahia, 145 séries em duas origens totalizam 4.350 ajustes; a seleção por perfil aumenta o atendimento em 3,89 pontos percentuais e o custo em 1,93\% frente à referência global. No piloto Walmart M5, 30 séries de dez lojas em três origens totalizam 1.350 ajustes. O preditivo por classe apresenta custo 13,9\% maior e atendimento 4,11 pontos inferior. Uma extensão exploratória por consequências previstas, com recuo à referência, reduz custo em 1,56\% e aumenta atendimento em 0,27 ponto; apenas três de 90 decisões mudam, e uma série concentra cerca de 92\% da economia. A metodologia permite distinguir ganhos da seleção, limitações do portfólio e restrições físicas do atendimento. As aplicações mostram compromissos distintos entre custo e serviço, sem estabelecer superioridade geral; a extensão desenvolvida no piloto exige confirmação independente.

\end{abstract}
\begin{keyword}
Inventory methodology \sep Modular decision support \sep Policy selection \sep Demand heterogeneity \sep Service level
\end{keyword}
\end{frontmatter}

\section{Introdução}
O gerente de logística e estoques precisa evitar custos desnecessários sem comprometer o atendimento. Selecionar políticas adequadas ao contexto da demanda é uma das decisões que podem apoiar esse objetivo. Vendas frequentes, intervalos sem demanda e picos ocasionais podem favorecer regras diferentes. A vantagem de contextualizar a escolha, porém, depende de reconhecer essas diferenças antes da decisão e convertê-las em resultados operacionais posteriores \cite{Axsater2006InventoryControl}.

A escolha da família de política constitui uma decisão distinta da calibração de seus parâmetros. Uma referência global pode usar a mesma família em toda a rede e ainda ajustar lotes e limiares ao histórico de cada série. A seleção de algoritmos fornece uma formulação para ir além dessa referência: características da instância orientam a escolha de um método de um portfólio \cite{rice1976,Kotthoff2014AlgorithmSelection}. No inventário, essa escolha precisa considerar tanto o custo quanto a viabilidade do atendimento.

A seleção contextual exige conectar três decisões: como representar a demanda, que evidência usar para recomendar e quando substituir a política de referência. Comparar candidatas retrospectivamente esclarece o potencial do portfólio; antecipar a melhor escolha requer informação disponível antes do teste. A avaliação também precisa identificar a origem de uma falha: erro do seletor, ausência de candidata viável ou impossibilidade física de atender ao alvo. Esse diagnóstico determina se a revisão deve incidir sobre a escolha, o portfólio ou as condições do sistema.

Propomos o \textit{Adaptive Inventory Policy Engine} (AIPE), uma metodologia modular para organizar essa decisão. O procedimento especifica o problema, representa o contexto, constrói uma matriz comparável de consequências, recomenda pelo \textit{Policy Selection Engine} (PSE) e avalia a escolha em dados reservados. O PSE admite associação por perfil ou estimação por série. Uma regra de intervenção explicita quando a candidata substitui a referência. Contratos de entrada, saída e informação permitem substituir métodos e reavaliar suas consequências mantendo o problema de decisão definido.

A contribuição do AIPE é formular essa sequência como um procedimento verificável: a evidência sustenta a recomendação, a regra de intervenção define sua aplicação e o diagnóstico identifica onde aperfeiçoá-la. Os controladores e seletores concretizam responsabilidades substituíveis dentro desse procedimento. Avaliamos tanto a compatibilidade entre essas implementações quanto sua utilidade incremental frente à referência global. A integração entre classificação e inventário possui antecedentes; o posicionamento do trabalho depende da articulação proposta e das decisões que ela permite examinar.

Examinamos o procedimento em vendas de uma rede varejista brasileira e em um piloto da base pública Walmart M5. A Bahia permite avaliar a seleção por perfil com vendas empresariais observadas sob condições simuladas; o M5 permite comparar implementações preditivas em um referencial público. As duas aplicações são analisadas separadamente, com consequências simuladas e calibrações identificadas. Os resultados revelam atendimento adicional com custo maior na Bahia, deterioração do preditivo por classe no M5 e melhoria concentrada quando a previsão de consequências é combinada com recuo à referência. Esse conjunto orienta a análise de quando intervir e de qual componente revisar.

\section{Trabalhos relacionados}
O intervalo médio entre demandas (\textit{average demand interval}, ADI) e o coeficiente de variação (CV) fundamentam a classificação ADI--CV$^2$, em que CV$^2$ é o quadrado do coeficiente. Essa classificação organiza padrões de frequência e tamanho da demanda \cite{syntetos2005}. Entretanto, classes estatísticas não determinam sozinhas a melhor reposição: custos, prazo, estoque inicial e atendimento também importam. Mohammaditabar et al.\ \cite{Mohammaditabar2012Integrated} integram classificação e escolha de políticas; Sheikh-Zadeh et al.\ \cite{SheikhZadeh2021PBIC} propõem classificação orientada ao desempenho em sistemas de múltiplos itens e escalões. O AIPE parte dessa relação entre classificação e desempenho e integra em seu procedimento uma regra explícita de intervenção e um diagnóstico das perdas frente ao portfólio e à viabilidade física. A avaliação deste artigo compara suas implementações internas; os métodos desses antecedentes permanecem fora da comparação empírica.

Na literatura da revista, Tadayonrad e Ndiaye \cite{sca_f81ad2b3e023} aproximam previsão e consequências operacionais. Theunissen et al.\ \cite{sca_c5c06dcbc9e5} discutem a informação decisória preservada pela classificação. Essa perspectiva orienta nossa avaliação pelo custo e pelo atendimento obtidos, em complemento à concordância entre famílias recomendadas. O método topológico discutido por esses autores oferece um antecedente conceitual, sem implementação nesta avaliação.

Zabraoui et al.\ \cite{sca_27e994c54848} comparam algoritmos de inventário e combinações de busca genética e aprendizado por reforço com dados M5. A combinação de algoritmo genético (\textit{genetic algorithm}, GA) e rede Q profunda (\textit{deep Q-network}, DQN), denominada GA--DQN, integra nosso portfólio, com seus componentes como controles, sem ser apresentada como contribuição algorítmica. Granularidade, filtros, custos e implementação local impedem uma reprodução exata de seu ranking.

Políticas limitadas de estoque-base oferecem outra referência para vendas perdidas \cite{Xin2021CappedBaseStock}. Suas propriedades teóricas não garantem desempenho para nossa calibração. Os resultados de reforço também dependem do estado, da recompensa e do orçamento \cite{Boute2022DRLInventoryRoadmap}. O AIPE separa essas implementações da escolha entre candidatas. A avaliação do seletor precisa, ainda, separar desenvolvimento e teste para evitar otimismo após escolher configurações \cite{CawleyTalbot2010}.

Francuz e Bányai \cite{FrancuzBanyai2026Selection} integram escolha de família e parâmetros em otimização multiobjetivo, usando aprendizado para aproximar o simulador. A aplicação tem demanda sintética sazonal e pedidos pendentes. Aprendizado orientado à consequência \cite{ElmachtoubGrigas2022SPO}, predição com rejeição \cite{GeifmanElYaniv2019SelectiveNet} e recuo à referência em aprendizado por reforço (\textit{reinforcement learning}, RL) \cite{Laroche2019SPIBB} também possuem antecedentes. O AIPE reúne essas ideias em um procedimento de seleção de famílias sob vendas perdidas, com referência calibrada e decomposição retrospectiva das falhas. A contribuição proposta é um procedimento verificável que registra a referência calibrada, a autorização para intervir e o diagnóstico das falhas de seleção, do portfólio e da viabilidade física. Classificação, predição de consequências e recuo fornecem os métodos constituintes. As aplicações examinam essa articulação na reposição com vendas perdidas; seus resultados não demonstram superioridade frente às integrações próximas, que permanecem sem confronto empírico neste estudo.

\section{AIPE: metodologia modular}
\subsection{Objeto e procedimento}
O AIPE apoia o gerente na escolha de uma família de reposição e de sua parametrização para um contexto de estoque, usando informação disponível antes da decisão. As séries loja--produto identificam contextos de demanda; diferenças entre elas não são interpretadas como medidas de eficiência dos gerentes ou das lojas. A aplicação começa pela definição de históricos e contexto, portfólio, ambiente de estoque, objetivo operacional e desenho de avaliação. A entrega pretendida reúne recomendação, parâmetros de aplicação, referência, alternativas, consequências esperadas e suporte da escolha, deixando a decisão de uso ao gerente. O registro distingue estimativas disponíveis na escolha e resultados conhecidos somente depois do horizonte de avaliação. Quando nenhuma candidata é viável, o registro conserva essa condição mesmo que uma contingência permita emitir a recomendação.

O procedimento tem cinco etapas: (1) especificar unidade, ambiente, objetivo e referência; (2) preparar históricos e representar o contexto; (3) calibrar candidatas e construir evidência de validação; (4) recomendar pelo PSE e aplicar a regra de intervenção; (5) avaliar em dados reservados e diagnosticar perdas. A Figura~\ref{fig:aipe} sintetiza a motivação; a Tabela~\ref{tab:modules} explicita os contratos metodológicos. Cada contrato define o que uma implementação recebe, o que entrega e quais informações pode usar.

\begin{figure}[htbp]\centering
\includegraphics[width=\linewidth]{figures/aipe_conceito_v2_ptbr.pdf}
\caption{AIPE: o contexto orienta a seleção de políticas. O PSE associa um perfil informado a uma família ou estima uma recomendação por série. Famílias A/B/C são ilustrativas; o benefício é avaliado operacionalmente.}\label{fig:aipe}
\end{figure}

\begin{table}[htbp]\centering\small
\caption{Contratos dos módulos do AIPE. Trocar o método preserva a responsabilidade e as restrições de informação.}\label{tab:modules}
\begin{tabular}{@{}p{2.4cm}p{4.0cm}p{4.8cm}@{}}\toprule
Módulo & Entrada e saída & Contrato e instanciação avaliada\\\midrule
Especificação & Requisitos $\to$ cenário e regra & Definir custo, serviço e referência; aqui, vendas perdidas e prazo fixo\\
Contexto & Histórico $\to$ características ou perfil & Usar informação anterior à escolha; descritores e regras fixas\\
Evidência & Cenário e candidatas $\to$ consequências & Ajustar no treino e medir na validação; sete candidatas clássicas, GA e RL\\
PSE & Contexto e evidência $\to$ recomendação & Preservar objetivo e corte temporal; perfil, classe vencedora ou consequências previstas\\
Avaliação & Recomendação fixada $\to$ comparação e diagnóstico & Teste reservado, referência competitiva e perdas identificáveis\\\bottomrule
\end{tabular}\end{table}

\subsection{Informação e configuração temporal}
A aplicação declara produto, local, datas e unidades, histórico de demanda ou vendas, custos, prazo e meta de atendimento. Saldo, compromissos e pedidos em trânsito são necessários para orientar a próxima compra real; na avaliação de políticas, podem ser substituídos por condições de cenário identificadas. Lotes mínimos, calendários e demais restrições aplicáveis delimitam as ações admissíveis. Cada entrada deve identificar origem observada, informada, estimada ou assumida em cenário, além de qualidade, transformações e disponibilidade temporal. Ausência de registro exige tratamento explícito.

O perfil pode ser fornecido ao PSE ou derivado do histórico admissível. Nas aplicações deste artigo, todos os perfis são derivados. A recomendação deve permitir consultar alternativas e rastrear a escolha até dados, regras e versões. Esse contrato é mais amplo que a persistência experimental atual, que ainda não reúne parâmetros finais de todas as políticas e trajetórias completas; os resultados não constituem ordens de compra executáveis.

A configuração distingue período básico de observação e ação, ciclo de análise, cadência de revisão do PSE, horizonte de planejamento, janela histórica e frequência de atualização dos modelos. A cadência informada para a empresa-fonte é uma reavaliação por ciclo comercial; essa é uma configuração de aplicação. Rever a recomendação pode conservar a família e os parâmetros, sem retreinar o seletor. Mudar a escala temporal exige compatibilizar agregação, incidência dos custos e prazo com a informação disponível. A flexibilidade proposta exige validação em cada configuração examinada.

\subsection{Substituição dos métodos e decisão de intervenção}
A modularidade permite avaliar uma alternativa em uma etapa mantendo definido o problema das demais. Um controlador recebe o estado disponível e produz um pedido válido; um descritor resume o histórico anterior à escolha; um seletor recomenda uma família do portfólio. A comparação preserva custos, definição do serviço, cronologia e unidades de análise. Quando uma alternativa exige outras ações ou informações, a adaptação de seu contrato precisa ser verificada antes da comparação.

Seja $\Pi$ o portfólio, $x_i$ a descrição do treino e $g_i$ o perfil. O PSE por perfil recomenda $m(g_i)$; o preditivo recomenda $f(x_i,g_i)$ ou $f(x_i)$. Ele pode aprender um rótulo ou estimar consequências por candidata. A referência global mantém a família, com parâmetros calibrados individualmente. A regra de intervenção escolhe entre candidata $\widehat\pi_i$ e referência $\pi_i^{\mathrm{ref}}$:
\[
 \pi_i^{\mathrm{AIPE}}=
 \begin{cases}\widehat\pi_i,&a_i=1,\\\pi_i^{\mathrm{ref}},&a_i=0.\end{cases}
\]
O indicador $a_i$ registra a decisão de intervir. Na implementação direta, toda recomendação é aplicada; na conservadora, a substituição exige o critério de benefício definido no desenvolvimento. Essa separação permite avaliar a candidata proposta e a regra que autoriza seu uso. Confiança de classe e serviço previsto são sinais para decidir, sujeitos a erro no atendimento realizado.

A avaliação exercita a substituição de controladores e seletores, mantendo a representação de contexto e o simulador. Nas séries avaliadas, todas as candidatas são calibradas exclusivamente no prefixo de treino. Suas trajetórias posteriores no teste permitem medir as recomendações e construir os oráculos retrospectivos. Essa despesa experimental difere do uso pretendido, em que evidências anteriores orientariam a família a calibrar; eventual economia computacional de implantação requer avaliação própria. A adaptação aqui estudada é a escolha por contexto entre origens; dentro de cada bloco de teste, a família permanece fixada por sete ciclos. As ações da política ao longo do bloco não constituem novas revisões do PSE.

A integridade da comparação exige uma matriz completa de candidatas, sementes e blocos. Imputação e codificação usam somente a partição de treinamento; recomendações, referências e intervenções são fixadas antes de medir o teste. A identificação de dados, configuração e código permite verificar quais resultados pertencem à mesma avaliação.

\subsection{Custo, serviço e escolha}
Para demanda $d_t$, quantidade atendida $v_t$, saldo $I_t$ e pedido $q_t$, avaliamos
\begin{equation}
 C=\sum_t[hI_t+p(d_t-v_t)+K\mathbf1(q_t>0)+cq_t],
 \qquad F=\frac{\sum_t v_t}{\sum_t d_t}.
 \label{eq:metrics}
\end{equation}
Os coeficientes $h$ e $p$ representam, respectivamente, custo de manutenção por unidade e ciclo e custo por unidade não atendida. O custo de emissão é $K$ por pedido, acrescido de $c$ por unidade pedida. O atendimento $F$ corresponde ao \textit{fill rate}, a fração da demanda atendida. Diferenças entre taxas de atendimento são expressas em pontos percentuais (p.p.). Com demanda zero, permanece indefinido; o custo continua contabilizado. Na instanciação, usamos $h=1$, $p=5$, $K=50$ e $c=0{,}5$, sem equivalência monetária calibrada.

As barras indicam médias aritméticas por caso série--origem, após agregar as sementes de cada candidata; serviços indefinidos são excluídos da média de atendimento. Na validação $A$, a regra global minimiza $\bar C_{\pi,A}$ entre candidatas com $\bar F_{\pi,A}\geq\alpha$, com $\alpha=0{,}70$. A meta média permite casos individualmente abaixo do alvo. O PSE por perfil aplica essa regra a cada grupo $A_g$. Casos sem demanda não entram na média de serviço. Sem candidata viável, escolhe-se maior serviço e menor custo, registrando contingência. Perfis sem suporte usam a global. Na Bahia, empate no custo viável prefere maior serviço e depois identificador; no M5, prefere identificador. Contingências empatadas em serviço preferem menor custo. Essas convenções são preservadas por base e registradas no código.

O rótulo supervisionado aplica a restrição à própria série de validação: menor custo viável ou, sem viável, maior serviço e menor custo. Se todas as opções têm serviço indefinido, utiliza menor custo e mantém a marcação de demanda zero. Classificador e seletores agregados não impõem a mesma distribuição da restrição. Incluímos a moda dos rótulos como controle para não atribuir toda diferença à capacidade preditiva.

\subsection{Características e perfis}
O preditivo recebe 17 descritores: ADI, CV$^2$, intermitência, proporção de zeros, maior e média das sequências de zeros, \textit{burstiness}, assimetria, curtose, entropia, tendência, autocorrelação, média, desvio-padrão, CV, extensão e número de ciclos positivos. O perfil categórico deriva desses descritores, sem constituir informação independente.

Os perfis usam regras ordenadas e fixas: esparso (ADI $\geq1{,}32$, CV$^2\geq0{,}49$, sequência de zeros $\geq5$); volatilidade/sazonalidade (autocorrelação $\geq0{,}25$ ou burstiness $\geq0{,}30$); tendência instável (tendência absoluta $\geq0{,}05$ ou entropia $\geq2$); baixo volume estável (média $\leq2$ e CV$^2<0{,}49$); e venda frequente (ADI $<1{,}32$). A primeira regra satisfeita define o perfil; o residual é esparso. ADI é $n/\max(n_+,1)$ e CV$^2$ inclui zeros, diferindo da taxonomia original baseada em tamanhos positivos. Os nomes são heurísticos.

Burstiness compara o desvio e a média dos intervalos positivos pela razão de sua diferença e soma, com regularização numérica. Tendência é a inclinação ajustada por mínimos quadrados ordinários dividida pela média; entropia usa dez classes de mesma largura. A autocorrelação na defasagem 17 recebe zero com até 18 ciclos ou variância desprezível. Essa convenção denota estimativa indisponível, não ausência de sazonalidade; com 23 ciclos, há somente seis pares.

\section{Desenho experimental}
A avaliação acompanha as decisões apoiadas pelo AIPE, sem medir adoção ou eficácia do uso por gerentes, em três perguntas. A seleção por perfil melhora o compromisso entre custo e atendimento frente à referência global? A estimação por série acrescenta valor e a regra de intervenção preserva esse valor? As falhas de atendimento decorrem da seleção, do portfólio ou de restrições físicas? A Bahia investiga a primeira pergunta; o piloto M5 permite examinar as três, incluindo a substituição do preditivo por classe por um seletor de consequências na mesma matriz de controladores.
\subsection{Bases e unidades de análise}
A Tabela~\ref{tab:datasets} distingue as avaliações. Na Bahia, transações complementares são somadas por loja, produto e ciclo; a preparação preenche com zero células sem agregado registrado. Esse preenchimento pressupõe cobertura das vendas e não comprova demanda nula. A coorte retrospectiva contém 145 séries, 115 lojas e nove produtos. Os 38 ciclos disponíveis, de 202301 a 202504, são unidades comerciais da fonte, não meses.

O M5 é uma fonte pública de vendas associada a uma competição de previsão \cite{Makridakis2022M5Accuracy}. Os códigos de lojas preservados do M5 combinam o estado e um identificador local: CA designa Califórnia, TX designa Texas e WI designa Wisconsin; por exemplo, CA\_2 identifica uma loja da Califórnia. O piloto sorteia 30 séries da categoria de alimentos (\texttt{FOODS}), três por loja nas dez lojas. Na primeira origem, o volume médio do treino define três estratos por loja, com uma escolha em cada estrato e semente 20261015. As séries são mantidas nas três origens. Usamos ciclos completos de 21 dias, sem remover picos por regras baseadas no intervalo interquartil.

\begin{table}[htbp]\centering\small
\caption{Matrizes completas analisadas. Casos série--origem e sementes não são séries independentes.}\label{tab:datasets}
\begin{tabular}{@{}lrr@{}}\toprule
Dimensão & Bahia & Piloto M5\\\midrule
Séries / lojas & 145 / 115 & 30 / 10\\
Origens de treino (ciclos) & 17 e 23 & 50, 64 e 78\\
Validação / teste por origem & 7 / 7 & 7 / 7\\
Casos série--origem & 290 & 90\\
Ajustes de políticas & 4.350 & 1.350\\
Testes sem demanda & 7 & 4\\
Seletores & Global e por perfil & Global, por perfil e preditivo\\\bottomrule
\end{tabular}\end{table}

A elegibilidade exige três ciclos positivos no prefixo. O sorteio M5 usa 885 séries elegíveis na primeira origem de uma coorte histórica de 2.408 identificadores. Não representa todas as FOODS nem séries elegíveis somente depois. O recorte é inspirado no contexto de Zabraoui et al., sem afirmar identidade com sua amostra. As conclusões quantitativas deste artigo se restringem ao piloto descrito.

\subsection{Simulador, portfólio e calibração}
O prazo exato é de dois ciclos. A política observa o estado anterior ao recebimento; entregas chegam antes da demanda. Quantidades podem ser fracionárias. Faltas são vendas perdidas, sem pedidos pendentes. Manutenção incide no saldo após demanda e o pedido é custeado na emissão, sem crédito terminal.

Validação e teste começam em estados comuns às candidatas, produzidos por uma referência de lote econômico de compra (\textit{economic order quantity}, EOQ): seis ciclos passados aquecem o sistema, e sua trajetória na validação determina estoque e pipeline do teste. Esses estados usam o passado, mas não representam a continuação particular de cada candidata. Escalas são calculadas no treino.

O portfólio contém EOQCycle, sSCycle, CappedBaseStock, GA, otimização proximal de políticas (\textit{proximal policy optimization}, PPO), GA--DQN e DQN. Denotamos por $\mu$ e $\sigma$ a média e o desvio-padrão da demanda no prefixo de treino. EOQCycle usa lote $\max(1,\sqrt{2\mu K/h})$ e ponto de pedido $2\mu+1{,}28\sigma\sqrt2$; sSCycle repõe até o ponto acrescido do lote. CappedBaseStock limita a quantidade da reposição até um alvo e calibra alvo e teto em grade $20\times12$. GA ajusta gatilho, lote e proteção com população 100 e 50 gerações. DQN usa Double DQN dueling; PPO usa otimização proximal. Ambos têm duas camadas de 64 unidades, 20 ações e no mínimo 20 mil transições, arredondadas ao episódio completo. GA--DQN acrescenta busca genética e mil demonstrações, mantendo gatilho e piso iguais nas demonstrações, treino e inferência.

Os quatro métodos estocásticos usam sementes 11, 29 e 47; as três referências determinísticas têm um ajuste. As sementes são promediadas antes da avaliação por caso. O híbrido tem esforço superior ao DQN, sem equivalência de orçamento.

Os objetivos de calibração incluem penalidades de atendimento, distintas do custo $C$ reportado. GA minimiza $C+30\max(0;0{,}70-F)\max(C;1)$ no treino. DQN, PPO e GA--DQN recebem recompensa de custo negativo e penalidade terminal $30\max(0;0{,}70-F)S$, com escala fixa $S=\max(p\sum_t d_t;1)$ calculada na demanda do episódio de treino.

A CappedBaseStock minimiza custo acrescido dessa penalidade com escala fixa. Na matriz Bahia reportada, acrescenta ainda $S[r(1-F)+\max(b-1;0)]$, em que $r$ é a fração de ciclos com falta e $b$ a razão regularizada entre variâncias de pedidos e demanda. Esses dois termos têm pesos unitários na Bahia e zero no M5. Penalidades orientam a calibração sem garantir o alvo; a seleção aplica depois a viabilidade observada na validação. Essa diferença impede atribuir contrastes entre bases somente à demanda. Cada matriz e código arquivado identificam a instanciação avaliada.

\subsection{Avaliação e desenvolvimento do piloto}
Em cada origem, uma loja inteira fica fora do conjunto que define a escolha. O PSE usa a validação das demais lojas e recomenda para o teste da loja reservada. A família global é comum às séries recomendadas nessa partição e origem; pode mudar entre lojas reservadas e entre origens. Controladores continuam calibrados no histórico da própria série: avaliamos transferência do seletor com histórico disponível, não implantação sem dados prévios.

No M5, XGBoost tem 200 árvores, profundidade três, taxa 0,05, sem subamostragem e com semente fixa. Imputação, codificação categórica e definição das classes são ajustadas dentro de cada partição de treinamento. Classes raras não são removidas; treinos de classe única usam classificador constante. Cada treinamento reúne 27 exemplos. A ablação sem perfil mantém os demais componentes. Não se usam origens posteriores para recomendar em uma origem anterior.

Uma implementação com recuo usa o PSE por perfil se o maior escore de classe for inferior a 0,60 ou a classe prevista tiver menos de três exemplos de treino. Não há busca de limiares no teste, mas a motivação veio dos resultados observados. Seu desempenho é exploratório, e o piloto passa a ser conjunto de desenvolvimento \cite{CawleyTalbot2010}. O escore não é probabilidade calibrada de viabilidade.

A comparação reduzida adicional reutiliza a matriz M5 e substitui a implementação do PSE. Duas florestas ExtraTrees estimam custo e serviço por candidata: 64 árvores, profundidade quatro, mínimo de três observações por folha e semente fixa. Custo usa $\log(1+C)$, transformado de volta; serviço exclui alvos indefinidos. A seleção direta escolhe o menor custo previsto viável; sem viável, prioriza serviço e depois custo.

A variante com recuo à global exige serviço previsto $\geq0{,}70$, não inferior ao previsto da referência, e economia acima de uma margem em $\{0;0{,}02;0{,}05;0{,}10;\infty\}$. A margem minimiza custo em validação interna por loja, sujeita a serviço médio não inferior à referência. Cada origem exclui uma loja externa; dentro das nove restantes, ajusta em oito e valida na nona. A margem é fixada antes da recomendação externa. Um classificador ExtraTrees com a mesma regularização controla a aprendizagem por rótulo, e uma ablação retira perfil do seletor por consequências sem recuo. Recomendações são gravadas antes da pontuação do teste, mas o piloto já orientou o desenvolvimento: a comparação permanece exploratória.

Na Bahia, os testes compartilham um ciclo; no M5, ocupam os ciclos 58--64, 72--78 e 86--92, sem sobreposição. Origens repetem séries. Intervalos exploratórios usam bootstrap pareado por loja, preservando séries e origens: 5.000 sorteios na Bahia, 2.000 na comparação de perfil e classe no M5 e 5.000 na comparação por consequências. Seletores e controladores permanecem fixos; a incerteza é condicional, sem novos treinamentos. Dez lojas do piloto limitam sua precisão.

\subsection{Oráculos e informação perfeita}
O melhor controlador único no teste é um diagnóstico retrospectivo. O oráculo do portfólio escolhe, por série--origem, a candidata de menor custo viável nas médias de teste entre sementes. Essas escolhas não fornecem rótulos ao treinamento. O acerto exige serviço mínimo e empate numérico no menor custo; concordância nominal é informada separadamente.

No M5, um programa linear inteiro misto (\textit{mixed-integer linear programming}, MILP) conhece toda a demanda de teste e otimiza pedidos sob os mesmos custos, prazo e condições iniciais. Quantidades são contínuas; variáveis binárias representam emissão e complementaridade entre saldo e falta. É um ótimo determinístico condicional, não um controlador estocástico implantável. Se 0,70 for fisicamente impossível antes das primeiras entregas, maximiza atendimento e depois minimiza custo. Soluções são reexecutadas no simulador e somente as comprovadamente ótimas recebem essa denominação.

A eficiência de custo é $100C^*/C_{\mathrm{PSE}}$. Aqui, $C^*$ denota o custo da referência retrospectiva considerada: o mínimo viável do portfólio ou o ótimo do MILP, conforme a comparação. Calculamos a razão apenas quando referência e recomendação satisfazem o mesmo critério de serviço; ambos os custos zero equivalem a 100\%. Serviço insuficiente não recebe crédito por custo baixo. A comparação com o MILP usa cada realização de treinamento viável: uma média entre sementes pode satisfazer o alvo sem que todas o satisfaçam. Repetir políticas determinísticas nessa análise não gera novas observações independentes.

\input{aipe_bahia_resultados}
A fronteira da Figura~\ref{fig:cost_service} situa o compromisso do portfólio, enquanto a Tabela~\ref{tab:ba2026_selection} mostra como a seleção por perfil modifica esse compromisso entre origens. A coincidência das recomendações na segunda origem delimita o período em que a segmentação acrescenta atendimento.
\input{m5_pse_20261015}
A Figura~\ref{fig:m5_comparison} permite relacionar as médias das estratégias à distribuição do atendimento por caso: o recuo aproxima o custo da referência, mas conserva falhas individuais. Esse contraste motiva examinar a previsão das consequências operacionais como outra implementação do PSE.
\input{aipe_consequencias_resultados}

\section{Discussão}
A relevância prática depende do contexto do gerente, da escala dos custos evitados e do atendimento preservado, sem exigir um percentual mínimo universal de economia. A avaliação deve fornecer magnitude, incerteza, deteriorações e rastreabilidade para esse julgamento. Ausência de diferença estatisticamente detectável não demonstra equivalência ou não inferioridade; uma garantia desse tipo exigiria margem e desenho próprios. Custos de cenário não autorizam converter diferenças em economia realizada em reais.

As duas aplicações mostram por que o AIPE avalia a seleção contextual frente a uma referência calibrada por série. A variação das vencedoras do portfólio oferece oportunidades de escolha, enquanto os resultados dos seletores revelam quanto desse potencial é efetivamente aproveitado. A compatibilidade entre módulos permite realizar essas comparações; o benefício operacional depende das consequências da recomendação em dados reservados.

Na Bahia, a seleção por perfil acrescenta atendimento com custo pontualmente maior; na segunda origem, coincide com a escolha global. Esse comportamento mostra que a utilidade da segmentação depende da janela e do critério operacional. A base brasileira fortalece a relevância aplicada da avaliação ao preservar heterogeneidade observada em vendas empresariais, enquanto a simulação e a possível censura por ruptura delimitam a interpretação econômica.

O piloto M5 evidencia a dificuldade de aprender a identidade da melhor família com pouco suporte. O XGBoost perde para a referência em custo e atendimento. A previsão de consequências preserva a magnitude operacional dos erros de escolha, mas sua implementação direta também perde serviço; o controle por classe com outra floresta altera os resultados. Portanto, a comparação não isola todo o efeito do alvo de aprendizagem.

A regra de intervenção torna explícito quanto o PSE modifica a referência e com que consequência. A variante por consequências conserva a global em 87/90 casos, melhora pontualmente custo e serviço e concentra cerca de 92\% da economia em uma série. A média favorável, portanto, depende de poucas decisões e deve ser lida junto à concentração do ganho e às deteriorações individuais. Esse resultado exploratório sustenta testar intervenções seletivas em dados externos ao desenvolvimento.

Os oráculos mostram que corrigir o seletor é apenas uma das intervenções possíveis. Havendo candidata viável, a escolha pode ser aprimorada; sem candidata viável, deve-se examinar o portfólio e sua calibração; diante de inviabilidade física, é necessário rever as condições do sistema. A distinção orienta qual componente do AIPE merece revisão e evita atribuir toda ruptura ao modelo preditivo.

O suporte do contexto também limita a transferência das recomendações. Na primeira origem M5, 826/885 séries elegíveis são esparsas, 57 têm tendência e duas pertencem ao grupo volatilidade/sazonalidade. Uma amostra maior pode manter essa concentração. A avaliação deve considerar o suporte de cada grupo e o objetivo de atendimento adotado: meta média, atendimento ponderado e garantia por série expressam exigências distintas. A comparação entre módulos só é interpretável quando conserva essa definição.

\section{Limitações e próximos testes}
Vendas podem omitir demanda censurada por falta. Coortes históricas, elegibilidade na primeira origem e horizontes de sete ciclos restringem população e período. Produtos compartilhados entre lojas podem gerar dependência não capturada inteiramente pelo bootstrap.

Não se trata de experimento causal de implantação nem de avaliação com usuários. A fotografia de estoque brasileiro disponível em maio de 2025 é posterior aos testes e não fornece seus saldos iniciais históricos. Custos não são financeiramente calibrados, não há suprimento conjunto e o estado inicial comum não reproduz trajetórias particulares. A calibração distinta da CappedBaseStock e os esforços desiguais limitam comparações entre algoritmos e bases.

As extensões com recuo foram desenvolvidas no piloto; suas métricas não confirmam independentemente a melhoria. A substituição foi exercitada nas candidatas e no PSE, sem ablação de todas as representações e ambientes. Assim, a portabilidade integral permanece uma hipótese de uso. A avaliação ampliada precisa examinar suporte por perfil e classe, restrições de serviço comparáveis e validação interna por loja para ajustes do seletor. Séries reservadas devem ficar fora dessas decisões. Harmonização da calibração e sensibilidades de custo, prazo e horizonte precedem conclusões mais amplas. Uma avaliação com revisão do PSE a cada ciclo deverá manter a trajetória própria de estoque e pedidos em trânsito de cada estratégia, inclusive ao mudar de família. Ela exige desenho sequencial adicional, pois os blocos atuais não validam essa rotina.

Uma avaliação ampliada, ainda fora dos resultados deste artigo, conserva as 145 séries brasileiras e os 2.408 identificadores históricos do M5. Nesse recorte, 885, 2.033 e 2.408 séries são elegíveis nas respectivas origens. O desenho usa pesos zero para ruptura e amplificação na calibração da CappedBaseStock em ambas as bases e cinco grupos internos de lojas na escolha da margem de recuo. Essas mudanças podem alterar as recomendações e exigem nova análise. Separar as séries do piloto das demais identifica o uso prévio dos dados, sem estabelecer independência por si só.

\section{Conclusão}
O AIPE propõe uma metodologia modular de apoio ao gerente de logística e estoques para examinar custos evitáveis e atendimento. Organiza a seleção contextual em cinco etapas, desde a especificação do problema até o diagnóstico dos resultados. A metodologia vincula a representação da série loja--produto à evidência de validação, à recomendação do PSE e à decisão de substituir a referência. Essa estrutura permite localizar tanto o benefício de uma implementação quanto o componente a revisar quando o atendimento falha.

A avaliação delimita onde essas escolhas ajudam. Na Bahia, o perfil acrescenta 3,89 pontos de atendimento com aumento de custo de 1,93\%. No piloto M5, a estimação por classe perde para a global; consequências com recuo reduzem custo em 1,56\% e acrescentam 0,27 ponto de atendimento, com apenas três intervenções e benefício concentrado. Assim, a evidência sustenta a execução e a utilidade diagnóstica do procedimento, enquanto seu benefício operacional permanece condicionado à instanciação.

A confirmação requer calibrações harmonizadas e avaliação da regra de intervenção em dados reservados, excluídos do desenvolvimento. O critério decisivo será a capacidade de reconhecer quando uma recomendação contextual compensa e quando preservar a referência, considerando serviço, custo e deteriorações.

\section*{Disponibilidade de dados e código}
Os dados brasileiros são proprietários e não são redistribuídos; o M5 é público. Código, configurações e fontes estão em \url{https://github.com/matheus013/msc}. Dados, bancos e modelos ficam fora da versão atual do repositório. Manifestos e snapshots locais identificam as matrizes; pacote arquivado e derivados autorizados deverão acompanhar a submissão. Não se promete acesso aos registros restritos mediante solicitação.

\section*{Autoria, financiamento e interesses}
Autoria, contribuições dos autores, financiamento, agradecimentos e interesses devem ser confirmados pelos autores responsáveis. O manuscrito em português requer tradução, revisão humana e fechamento das declarações antes da submissão.

\section*{Uso de inteligência artificial na preparação do manuscrito}
OpenAI Codex auxiliou organização bibliográfica, leitura de código, análise, visualização e redação. A figura conceitual usa desenhos vetoriais programáticos; os gráficos quantitativos derivam das matrizes identificadas. Conferência e responsabilidade pelo texto final cabem aos autores humanos. Sistemas de inteligência artificial não são autores.

\bibliographystyle{elsarticle-num}
\bibliography{references_ptbr}
\end{document}
